\section{Conclusion}
\label{sec:conclusion}

We have presented \textsc{FinOPD}, a financial multi-agent framework that breaks the evolution deadlock afflicting current systems by introducing on-policy self-distillation as the mechanism for experience-driven parametric improvement. The same model serves as both student and teacher through differential conditioning on hindsight risk-adjusted returns, with Shapley credit assignment disentangling contributions across five specialist agents and per-token KL clipping stabilizing training in the presence of stylistic token dominance. A complementary non-parametric track consolidates winning geometry-factor-action triples into a retrieval-augmented belief memory, providing rapid access to proven patterns without the dilution inherent in parametric averaging.

The evolution curve, our central empirical contribution, demonstrates that policy-level self-evolution in financial agents is both achievable and measurable: Sharpe ratio improves monotonically across self-distillation iterations under strict post-cutoff evaluation, ruling out pretraining leakage as an explanation. Across Dow-30 and CSI-300 benchmarks, \textsc{FinOPD} achieves consistent improvements over twelve baselines spanning numerical models, LLM multi-agent systems, and factor mining approaches.

Several limitations warrant discussion. The stability of token-level JSD matching under the high-variance reward signals characteristic of financial trajectories requires careful hyperparameter tuning; our alignment guard and adversarial refresh mitigate but do not eliminate drift risk. The belief store's capacity constraints introduce a forgetting mechanism whose optimal parameterization remains an open question. The factor library, while substantially larger than prior work at 160 candidates, remains modest compared to industrial-scale factor pools; extending the router to dynamically incorporate newly mined factors is a natural direction for future work.

More broadly, \textsc{FinOPD} establishes that the on-policy self-distillation paradigm, originally developed for single-model mathematical reasoning, transfers effectively to multi-agent financial decision-making when augmented with appropriate credit assignment and memory consolidation mechanisms. We release the codebase, the FinChartGeometry-50K dataset, and all trained checkpoints to facilitate reproducibility and further research.
